\section{Data Pipeline: Centroid Companion Files}
\label{sec:data-pipeline}

\subsection{File Format}

Phase~2 \cvdgeo{} requires a centroid companion file alongside the adjacency
file.
The file is named \texttt{\{state\}\_centroids\_\{year\}.json} and lives in
the same directory as the existing adjacency and geoids files:

\begin{verbatim}
data/2020/
  north_carolina_adjacency_2020.json
  north_carolina_geoids_2020.json
  north_carolina_centroids_2020.json   <-- Phase 2 addition
\end{verbatim}

The JSON structure is a single array of \texttt{[longitude, latitude]} pairs
in WGS84, one entry per census tract in the same index order as the
adjacency file:

\begin{verbatim}
{
  "centroids": [
    [-78.8740, 35.2271],   // tract 0
    [-80.8431, 35.2270],   // tract 1
    ...                    // one entry per tract
  ]
}
\end{verbatim}

\paragraph{Index alignment.}
The centroid file must have exactly $m$ entries, where $m$ is the number of
tracts in the adjacency file.
If $\mathrm{len}(\mathrm{centroids}) \ne m$, the loader returns an error:

\begin{verbatim}
ERROR: centroid file length mismatch: got 2659, expected 2660
       Check that centroids file was generated from the same
       adjacency file as the graph.
\end{verbatim}

\paragraph{Centroid definition.}
Each entry is the \emph{population-weighted polygon centroid} of the
census-tract boundary: the interior point $(\bar{\lambda}, \bar{\phi})$
minimising the average squared distance to the tract polygon boundary,
weighted by population distribution within the tract.
For most census tracts this is close to the geometric centroid of the polygon.
It is derived from the TIGER/Line shapefile polygon at fetch time
(Section~\ref{subsec:fetch}).

\subsection{Fetching Centroids}
\label{subsec:fetch}

Centroid companion files are generated by:

\begin{verbatim}
redist fetch --type centroids --states NC --year 2020
\end{verbatim}

This command computes population-weighted polygon centroids from the
TIGER/Line shapefile polygons already downloaded as part of the standard
\texttt{redist fetch} pipeline and writes them to the adjacency directory.
To fetch centroids for all 50 states:

\begin{verbatim}
redist fetch --type centroids --year 2020 --workers 8
\end{verbatim}

The fetch step is idempotent: if the centroid file already exists and its
length matches the adjacency file, the fetch is skipped.

\paragraph{CLI error when centroids absent.}
If \texttt{--cvd-metric geographic} is specified and the centroid companion
file is absent, the CLI returns:

\begin{verbatim}
ERROR: --cvd-metric geographic requires tract centroid data.
       Run: redist fetch --type centroids --states NC --year 2020
\end{verbatim}

No fallback to Phase~1 occurs silently: if the user requests geographic
metric and centroids are unavailable, the run fails with a clear diagnostic.

\subsection{LoadedGraph Integration}

The \texttt{LoadedGraph} struct gains one new field:

\begin{verbatim}
pub struct LoadedGraph {
    // ... existing fields ...
    /// Population-weighted centroid (WGS84 lon/lat) for each tract polygon.
    /// Empty if centroid companion file is absent (Phase 1 fallback; no error).
    pub tract_centroids: Vec<(f64, f64)>,
}
\end{verbatim}

Loading behaviour:
\begin{itemize}
\item After loading the adjacency file, look for
  \texttt{\{stem\}\_centroids\_\{year\}.json} in the same directory.
\item If present and length-valid: load into \texttt{tract\_centroids}.
\item If absent: leave \texttt{tract\_centroids} empty (no error at load
  time; error is deferred to Phase~2 dispatch if \texttt{--cvd-metric geographic}
  is requested).
\item If present but length-mismatched: return error immediately at load time.
\end{itemize}

This design preserves full backward compatibility: all existing Phase~1
\cvd{} runs, all \metis{} runs, and all other structure-layer algorithms
are unaffected by the presence or absence of the centroid field.

\subsection{Consistency with redist-map}

The Albers projection constants in \texttt{albers\_project}
(Definition~\ref{def:albers}) must match the constants used in
\texttt{redist-map} for area and Polsby-Popper computation.
Both use EPSG:5070 with central meridian $-96^{\circ}$, parallels
$29.5^{\circ}/45.5^{\circ}$, and $R = 6{,}371{,}000$\,m.
A platform invariant test asserts that
$\mathrm{proj}(-96.0^{\circ}, 37.5^{\circ})$ returns $(|x| < 1{,}000\,\mathrm{m}, y > 0)$:
the central meridian projects to near-zero $x$ (within 1\,km) and positive $y$
(positive northing in the projection).
